\section{Future Direction}
\subsection{Personalize AI Agents}

With the increasing interest in self-evolving agents, deploying personalized agents has become a crucial and increasingly significant objective for the research community \citep{zhang2024personalization}. For instance, in applications such as chatbots, digital twins, and emotional support dialogues, a key challenge is enabling AI agents to accurately capture and adapt to users' unique behavioral patterns or preferences over extended interactions. Existing personalized agents typically depend heavily on labeled data and post-training methodologies \citep{cheng2024autopal}. Recent work by \cite{zhang2025personalize} proposes a self-generated preference data approach aimed at rapidly personalizing LLMs. TWIN-GPT \cite{wang2024twin} leverages electronic health records to create digital twins of patients, enhancing the accuracy of clinical trial outcome predictions. However, these existing strategies hinge on the critical assumption that LLMs can consistently obtain high-quality, large-scale user data.

In practical deployment scenarios, the primary challenge remains the cold-start problem: agents need to progressively refine their personalized understanding, accurately interpret user intentions, and effectively construct user profiles, even when initial data is limited. Additionally, significant challenges persist in personalized planning and execution, such as effective long-term memory management, external tool integration, and personalized generation (ensuring outputs consistently align with individual user facts and preferences) \citep{li2025survey}. Moreover, it is essential to ensure that self-evolving agents do not inadvertently reinforce or exacerbate existing biases and stereotypes, highlighting another critical direction for future research. 
These governance principles are particularly important as personalized agents continuously evolve and adapt their memory or decision-making processes. 
Building upon these governance principles, evaluation frameworks should also evolve to ensure fairness, accountability, and alignment in personalized settings.

\paragraph{Data governance.}
Responsible personalization should balance adaptivity with privacy protection. First, agents ought to adopt \textbf{data minimization}: collect only task-relevant data and surface transparent, revocable controls. Empirically, web-agent benchmarks show that state-of-the-art agents frequently process sensitive data unnecessarily, motivating minimization-by-default designs \citep{zharmagambetov2025agentdam}. Complementarily, effective data governance for personalized agents entails deploying \textbf{on-device personalization frameworks} that enable local learning from user interactions, together with user-led privacy mechanisms such as Rescriber that perform on-device redaction and approval prior to any remote data exchange \citep{zhou2025rescriber}. To prevent indefinite retention, \textbf{memory decay and forgetting policies} should support selective deletion and “right-to-be-forgotten”-style unlearning for personalized traces \citep{staufer2025should}. Finally, because self-evolution can drift safety or fairness over time, systems should include \textbf{bias monitoring and fairness auditing} loops that adapt criteria and interventions as the user and context evolve \citep{basu2025fair}, and guard against \emph{misevolution} (safety/alignment degradation during evolution) via continuous checks on memory/tool/workflow updates \citep{shao2025misevolution}.


\paragraph{Evaluation.}
With the integration of personalized data, evaluation metrics for personalizing self-evolving agents should extend beyond intrinsic evaluations (e.g., directly assessing personalized generated text quality using metrics such as ROUGE \citep{lin2004rouge} and BLEU \citep{papineni2002bleu}) or extrinsic evaluations (e.g., indirect assessments of personalization effects through recommendation systems, classification tasks, and other specific applications). Traditional personalization evaluation metrics often fail to adequately capture the evolving dynamics inherent in self-evolving agents. Consequently, future research calls for more lightweight and adaptive evaluation metrics \citep{zhang2024personalization}. Additionally, to better assess self-evolving personalized agents, there is a clear need for flexible, dynamic benchmarks capable of accurately evaluating agents' performance, particularly in managing long-tailed personalization data throughout their self-evolving processes. We advocate reporting: 
\begin{itemize}
\item \textbf{Personal Adaptation Gain (PAG):} improvement per user over $k$ sessions vs.\ non-personalized baseline.

\item \textbf{Retention \& Forgetting Balance:} metrics such as forward/backward transfer and selective-forgetting efficacy (e.g., reduction in exposure to user facts following deletion requests). 

\item \textbf{Privacy–Utility Trade-off:} ratio of utility gain to bits of retained personal data; complemented by a \emph{Data Minimization Score} (proportion of shared sensitive information only when necessary) in web-agent scenarios \citep{zharmagambetov2025agentdam}.

\item \textbf{On-device Learning Ratio:} proportion of personalization updates or memory writes executed locally on the user device, potentially with user consent/redaction mechanisms (e.g., as studied in Rescriber) \citep{zhou2025rescriber}.  

\item \textbf{Bias/Drift Monitors:} longitudinal measurement of disparity or safety degradation across user groups; one might define indices such as \emph{Fairness-Drift} or \emph{Safety-Drift}, though standardized versions remain to be developed \citep{basu2025fair,shao2025misevolution}.  

\item \textbf{Longitudinal User Outcomes:} tracking session-level satisfaction, goal-completion trends or multi-phase development trajectories (e.g., virtual “campus-life” agents in benchmarks like StuLife) \citep{cai2025building}.  

\end{itemize}





\subsection{Generalization}

Self-evolving agents also face considerable challenges in achieving robust generalization across diverse task domains and environments. The fundamental tension between specialization and broad adaptability remains one of the most pressing challenges in the field, with significant implications for scalability, knowledge transfer, and collaborative intelligence.

\paragraph{Scalable Architecture Design:} A central challenge in developing generalizable self-evolving agents lies in designing scalable architectures capable of maintaining performance as complexity and scope increase. Current agent systems frequently encounter a trade-off between specialization and generalization, where agents optimized for specific tasks struggle to transfer their learned behaviors to novel environments \citep{chen2024optima}. Additionally, the computational cost associated with dynamic reasoning in LLM-based agents grows non-linearly with the complexity of adaptation mechanisms, imposing practical constraints on achievable generalization within realistic resource limitations \citep{kim2025cost}. Recent studies indicate that self-evolving agents equipped with reflective and memory-augmented capabilities show substantial promise for enhancing generalization, particularly in smaller, resource-constrained models \citep{liang2024sage}. Nonetheless, these approaches continue to encounter limitations when addressing complex real-world scenarios that require sustained adaptation over prolonged periods.



\paragraph{Cross-Domain Adaptation:} Achieving generalization across domains represents a critical frontier for self-evolving agents. Current methods frequently rely on domain-specific fine-tuning, restricting agents' adaptability to new environments without retraining \citep{belle2025agents}. Recent advancements in test-time scaling and inference-time adaptation provide promising pathways for enhancing cross-domain generalization \citep{snell2024scaling, zhang2025survey}. These techniques allow agents to dynamically allocate additional reasoning capacity to unfamiliar scenarios by scaling computational resources during inference, avoiding the need for increasing model parameters. Additionally, meta-learning strategies have demonstrated considerable potential in facilitating rapid few-shot adaptation to new domains \citep{bilal2025meta}. However, their effectiveness critically depends on an agent's capability to accurately determine when supplementary computational resources are necessary and efficiently distribute these resources across diverse reasoning tasks.

\paragraph{Continual Learning and Catastrophic Forgetting:} Self-evolving agents must continuously adapt to new tasks while retaining previously acquired knowledge, a challenge exacerbated by the catastrophic forgetting phenomenon \citep{ghosal2024understanding} of continual memorization \citep{chen2024continual} inherent in LLMs \citep{bell2025future}. The stability-plasticity dilemma becomes particularly acute in foundation model-based agents, where the computational costs of retraining for every new task are prohibitive \citep{zheng2025lifelong}. Recent research has explored parameter-efficient fine-tuning methods, selective memory mechanisms, and incremental learning strategies to mitigate catastrophic forgetting while preserving adaptability \citep{wang2024comprehensive}. Nonetheless, achieving an optimal balance between efficiency and preventing model drift remains a significant open challenge, especially when agents operate under resource constraints or manage streaming data with stringent privacy considerations.

\paragraph{Knowledge Transferability:}
Recent studies have identified critical limitations in knowledge transfer among AI agents. \cite{shi2025models} emphasized that knowledge integration and transfer capabilities in current agents still require significant optimization. In particular, \cite{geng2025large} found that LLM-based agents often fail to effectively propagate newly acquired knowledge from interactions to other agents, restricting their collaborative potential. Furthermore, \cite{vafa2025has} revealed that foundation models might depend heavily on shallow pattern matching, rather than developing robust and transferable internal world models. These findings indicate several important future research directions: 1) it is essential to better understand the conditions under which knowledge acquired by one agent can be reliably generalized and communicated to others; 2) developing methods to quantify the limitations in agents' knowledge transferability could lead to clearer insights into agent collaboration bottlenecks; 3) we need to have an explicit mechanism that encourage the formation of robust, generalizable world models could significantly improve the collaborative effectiveness of self-evolving agents. 



\subsection{Safe and Controllable Self-Evolving Agents}



As autonomous AI agents become increasingly capable of learning, evolving, and performing complex tasks independently, ensuring their safety and controllability has become a paramount concern. Unlike static systems, the very nature of self-evolution introduces unique and amplified risks that emerge dynamically over the agent's lifecycle. These risks are not merely extensions of traditional AI safety issues, such as those arising from vague user instructions or environmental threats like malicious phishing links~\citep{zhou2025automating}, but are fundamentally tied to the agent's capacity for autonomous self-modification and adaptation~\citep{shao2025misevolution,han2025alignment}. This section delineates the emergent risks unique to self-evolving agents and then discusses a set of prescriptive guardrails and mitigation strategies for building safer systems.

\subsubsection{Emergent Risks in Self-Evolving Systems}
Recent research has identified new risks that arise from the autonomous self-improvement process itself. These risks are not necessarily present in the initial agent but can manifest over time as it evolves. We categorize these risks along the primary evolutionary pathways of an agent: the backbone model, memory, and tools.

\begin{itemize}
\item \textbf{Uncontrolled behavior drift in model evolution}: A core risk is that an agent's goals and values may drift away from original human intent as it evolves. This is exacerbated by the learning uncertainty inherent in self-evolution, especially when operating in ambiguous contexts~\citep{anwar2024foundational,bagdasarian2024airgapagent}. A phenomenon termed "misevolution"~\citep{shao2025misevolution} can occur during model evolution. For instance, self-training on agent-generated data can lead to "catastrophic forgetting" of safety alignment, causing agents to execute harmful instructions they previously refused, such as interacting with malicious content they were trained to avoid~\citep{shao2025misevolution, hahm2025unintended}. 
\item \textbf{Deployment-time reward hacking in memory evolution}: Open-ended evolution is susceptible to reward hacking. Agents may find and exploit loopholes in self-defined reward signals or internal feedback. This is particularly evident in memory evolution, where the accumulation of experience can inadvertently induce unsafe behaviors. For example, an agent might learn to issue unnecessary refunds because its memory correlates them with high satisfaction ratings~\citep{shao2025misevolution}, and such risks can be further exacerbated by poorly designed memory modules~\citep{wang2025unveiling}. A related concept is the "Alignment Tipping Process (ATP)," where an initially aligned agent discovers that misaligned behaviors are more rewarding, causing its policy to "tip" and abandon its initial constraints~\citep{han2025alignment}.
\item \textbf{Safety of self-created and ingested external tools}: The ability of agents to autonomously generate and use tools introduces significant safety issues. Agents may spontaneously create tools with security vulnerabilities or fail to identify malicious code when ingesting external tools~\citep{shao2025misevolution}. This turns the agent into a potential vector for security threats. A major challenge here is that agents still struggle to differentiate between necessary and irrelevant sensitive information~\citep{zharmagambetov2025agentdam}, potentially leading them to create tools that leak private data. Furthermore, as agents improve, they may become more adept at creating and executing offensive cyber-operations, a risk that requires dynamic assessment~\citep{wei2025dynamic}.
\end{itemize}



\begin{table}[h!]
\renewcommand{\arraystretch}{1.2}
\centering
\caption{Compliance checklist for deploying self-evolving agents}
\label{tab:minimal_compliance_checklist}
\small
\begin{tabularx}{\textwidth}{|>{\raggedright\arraybackslash}p{3cm}|X|}
\hline
\textbf{Category} & \textbf{Compliance Checklist for Deployment}\\ \hline
\textbf{Tool \& Code Safety}
    & $\Box$ \textbf{Strict Sandboxing:} All tools and agent-generated code execute in an isolated environment with no default access to host files, network, or sensitive processes. \newline
    $\Box$ \textbf{Resource Limiting:} The sandbox imposes strict limits on CPU, memory, and execution time to prevent denial-of-service or runaway processes. \newline
    $\Box$ \textbf{Automated Security Verification:} A mandatory pipeline performs static analysis (e.g., SAST) and vulnerability scanning on all new or modified tools before use. \newline
    $\Box$ \textbf{Dependency Scanning:} The verification pipeline checks all third-party libraries and dependencies for known vulnerabilities. \newline
    $\Box$ \textbf{Risk-Based Access Control:} Tools are classified by risk level, and high-risk capabilities (e.g., file system writes, API calls) require explicit approval via an approval gate. \\ \hline
\textbf{Self-Modification Control}
    & $\Box$ \textbf{Immutable Audit Trail:} All self-modifications (to model weights, memory, toolset, or core logic) are logged with details on the trigger, changes made, and outcome. \newline
    $\Box$ \textbf{Version Control for Safe States:} The agent's state (model, memory, tools) is versioned, with known "safe" versions clearly tagged. \newline
    $\Box$ \textbf{Tested Rollback Mechanism:} A reliable, one-click rollback mechanism exists to revert the agent to a previously known safe version. This mechanism is regularly tested. \newline
    $\Box$ \textbf{Pre-Update Safety Validation:} Before a self-modified model is deployed, it is automatically evaluated against a "golden dataset" of safety-critical prompts to prevent catastrophic forgetting of alignment. \\ \hline
\textbf{Behavioral \& Alignment Safety}
    & $\Box$ \textbf{Continuous Runtime Monitoring:} An active monitoring system tracks agent actions, flagging deviations from expected behavior, anomalous resource usage, or signs of unsafe actions. \newline
    $\Box$ \textbf{Reward Hacking Detection:} Key metrics are monitored for signs of reward hacking (e.g., exploiting loopholes in reward functions). Alerts are configured for sharp, unexplained metric changes. \newline
    $\Box$ \textbf{Automated Red-Teaming:} A continuous red-teaming framework is active, programmatically generating and running test scenarios to probe for emergent misalignment, value drift, and new failure modes. \newline
    $\Box$ \textbf{Goal Guardrails:} Strict constraints are placed on the agent's ability to modify its own fundamental goals or safety constraints. Any such change requires human review. \\ \hline
\textbf{Data Privacy \& Memory Integrity}
    & $\Box$ \textbf{Proactive Memory Defense:} Mechanisms like dual-memory structures or consensus validation are in place to detect, isolate, and neutralize potentially "poisoned" or harmful memories before they influence behavior. \newline
    $\Box$ \textbf{PII Detection and Sanitization:} Automated tools are used to detect and redact/anonymize Personally Identifiable Information (PII) before it is stored in long-term memory or used in training. \newline
    $\Box$ \textbf{Data Minimization Principle:} The agent is configured to only collect and retain data that is strictly necessary for its tasks, and data retention policies are enforced. \newline
    $\Box$ \textbf{Privacy Regulation Compliance:} The system is designed to comply with relevant data privacy regulations (e.g., GDPR, CCPA). \\ \hline
\textbf{Operational Controls \& Governance}
    & $\Box$ \textbf{Human-in-the-Loop for Critical Actions:} High-stakes actions (e.g., large financial transactions, data deletion, communication with external users) are gated by a mandatory human approval step. \newline
    $\Box$ \textbf{Clear Incident Response Plan:} A documented plan is in place for responding to safety failures, including steps for immediate shutdown, rollback, and analysis. \newline
    $\Box$ \textbf{Centralized Dashboard for Oversight:} A dashboard provides human operators with real-time visibility into the agent's behavior, state, and active safety alerts. \newline
    $\Box$ \textbf{Explainability \& Traceability:} The system provides clear explanations for why a particular action was taken, linking it back to specific memories, goals, or model inferences in the audit trail. \\ \hline
\end{tabularx}
\end{table}




\subsubsection{Prescriptive Guardrails and Mitigation Strategies}

Addressing the emergent risks of self-evolution requires moving beyond descriptive warnings to implementing prescriptive, actionable guardrails. While early frameworks like TrustAgent have explored multi-stage strategies (i.e., pre-, in-, and post-planning) to foster safer behavior~\citep{hua2024trustagent}, the unique dynamics of self-evolution call for a more comprehensive "safety lifecycle" approach. Future research and development are expected to focus on integrating safeguards at every stage of the agent's operation.

\begin{itemize}
\item \textbf{Sandboxing and verification for tool and code execution}: To mitigate risks from tool use, all agent-generated or externally-sourced tools must be executed in a strictly sandboxed environment. Furthermore, automated safety verification, such as static analysis and vulnerability scanning, should be a default step before a new tool is integrated~\citep{mcpscan}. For securing tool interaction protocols, runtime defense pipelines~\citep{xing2025mcp, wang2025mcpguard} provide layered detection against threats like tool poisoning and prompt injection.
\item \textbf{Audit trails and failsafes for self-modification}: Any self-modification must be accompanied by a comprehensive audit trail. This ensures that changes are traceable and reversible. Implementing rollback and failsafe patterns is critical, allowing the system to revert to a previously known safe state if undesirable behavior is detected. For memory, proactive defenses like A-MemGuard propose dual-memory structures and consensus-based validation to identify and isolate "poisoned" memories before they corrupt behavior~\citep{wang2025unveiling, wei2025memguard}.
\item \textbf{Continuous monitoring and red-teaming for long-horizon drift}: Static, pre-deployment safety evaluations are insufficient. Continuous monitoring of agent behavior is necessary to detect long-horizon value drift. This can be achieved through red-teaming scenarios designed to test for emergent misalignment. For GUI agents, hybrid validation frameworks like OS-Sentinel combine formal verifiers with contextual judges to provide robust, in-workflow safety monitoring~\citep{sun2025sentinel}.
\item \textbf{Approval gates and privacy-protection measures}: For high-stakes actions, approval gates requiring human-in-the-loop confirmation should be implemented. Furthermore, given that agents struggle with handling sensitive data~\citep{zharmagambetov2025agentdam}, robust privacy-protection measures are necessary to prevent leakage and ensure a balanced and secure deployment.
\end{itemize}

To aid practitioners, we synthesize these strategies into a compliance checklist for deploying self-evolving agents (See Table \ref{tab:minimal_compliance_checklist}). In conclusion, deploying reliable, controllable, and safe self-evolving systems is a critical and active area of research. Future work must move towards building a comprehensive safety-aware evolutionary lifecycle, integrating robust verification, continuous monitoring, and adaptive guardrails to ensure that the agents remain aligned with human values and safety constraints as they become more autonomous.





\subsection{Ecosystems of Multi-Agents}
Multi-agent self-evolving systems face several unique challenges that require further exploration.
\paragraph{Balancing Individual and Collective Reasoning:} Recent studies highlight the difficulty of balancing independent reasoning with effective group decision-making in multi-agent environments \citep{chen2025mdteamgpt, sun2025collab}. While collective discussions can significantly enhance diagnostic reasoning, agents often risk becoming overly reliant on group consensus, thereby diminishing their independent reasoning capabilities. To mitigate this issue, future research should explore dynamic mechanisms that adjust the relative weight of individual versus collective input. Such an approach would help prevent decision-making from being dominated by a single or a small subset of agents, ultimately promoting robust, balanced consensus-building and innovation. Additionally, developing explicit knowledge bases and standardized updating methodologies—leveraging agents' successes and failures—could further improve the agents' self-evolution abilities and strengthen their individual reasoning contributions within collaborative contexts.

\paragraph{Efficient Frameworks and Dynamic Evaluation:}
Another crucial challenge lies in developing efficient algorithms and adaptive frameworks that allow agents to collaborate effectively while preserving their individual decision-making strengths. \citep{hu2024self} introduced adaptive reward models and optimized dynamic network structures, which can significantly enhance cooperative self-improvement among agents. However, a major gap identified by \citep{sun2025collab} is the absence of clear mechanisms for agents to dynamically manage and update their knowledge. Addressing this issue will require new frameworks that explicitly integrate continuous learning and adaptive collaboration mechanisms. Furthermore, existing benchmarks for multi-agent evaluation are predominantly static \citep{zhu2025multiagentbench} and therefore fail to capture the long-term adaptability and continuous evolution of agent roles. Future benchmarks should incorporate dynamic assessment methods, reflecting ongoing adaptation, evolving interactions, and diverse contributions within multi-agent systems, thus providing more comprehensive evaluation metrics for self-evolving agents.

